\section{La apuesta de largo plazo de CUDA: cómo un activo técnico se convierte en un activo estratégico}

\subsection{El costo de poner CUDA en GeForce}

Uno de los fragmentos históricos más valiosos de la entrevista es cuando Jensen recuerda ``putting CUDA on GeForce''. Dice explícitamente que en su momento fue una decisión que ``could not afford to do'': la presión financiera de corto plazo era grande, pero a largo plazo era una medida necesaria. La razón es directa: sin una base instalada no hay ecosistema de desarrolladores; sin ecosistema, la ventaja de la plataforma no es sostenible.

\begin{importantbox}{La base instalada es la variable estratégica de primer orden}
Jensen atribuye directamente el foso defensivo de CUDA a la base instalada. No es la ventaja de una generación de bibliotecas de operadores, ni la ventaja en un benchmark determinado, sino la variable compuesta de ``mentalidad del desarrollador + acumulación de software + velocidad de iteración''.
\end{importantbox}

\begin{figure}[H]
\centering
\includegraphics[width=0.92\textwidth]{frame-002.jpg}
\caption{Parte media de la entrevista: repaso de la decisión histórica clave entre CUDA y GeForce\protect\footnotemark}
\end{figure}
\footnotetext{Intervalo de tiempo en el video: 00:45:10--00:45:28. La imagen corresponde al momento en que Jensen habla de la expansión del ecosistema y la ruta de sistemas.}

\subsection{Por qué OpenCL no formó un volante equivalente}

El planteamiento de Jensen no niega la viabilidad de tecnologías alternativas, sino que enfatiza que la construcción de un ecosistema es una función del tiempo. Incluso si existe una alternativa a nivel de interfaz, la cadena de herramientas para desarrolladores, la experiencia de ajuste, la adaptación a los marcos de trabajo y la inercia de despliegue generan un costo de migración alto. Esto abre una brecha enorme entre ``sustituible a nivel técnico'' y ``sustituible a nivel comercial''.

\begin{warningbox}{Una ilusión en la competencia entre plataformas}
Muchos equipos confunden ``compatibilidad de sintaxis'' con ``compatibilidad de ecosistema''. En realidad, el costo central de migrar de ecosistema proviene del ajuste de rendimiento, el manejo de fallas y la experiencia de ingeniería acumulada, aspectos que rara vez aparecen en las tablas comparativas de API.
\end{warningbox}

\subsection{Los nuevos requisitos de CUDA al entrar en la era de la IA}

En la segunda mitad de la entrevista se menciona que las versiones de CUDA siguen iterando (por ejemplo, la generación 13.x) y se enfatiza que su valor no está solo en el training, sino también en inference, en el pipeline de agentes y en la colaboración entre múltiples modelos, entre otros escenarios nuevos. Si la plataforma quiere seguir liderando, necesita mantener una tensión entre ``generalidad'' y ``capacidad especializada'': debe poder soportar nuevas cargas de trabajo sin perder densidad de rendimiento.

\begin{knowledgebox}{De ``saber usar CUDA'' a ``saber operar CUDA''}
Del lado empresarial, el verdadero umbral de CUDA se está desplazando hacia la capacidad de operación de sistemas: cómo hacer aislamiento multiinquilino, escalado elástico de inference, programación de clústeres heterogéneos, degradación ante fallas y gobernanza de costos. La capacidad de desarrollo y la capacidad de operación de plataformas se están fusionando.
\end{knowledgebox}

\subsection{Resumen de la sección}

El valor de largo plazo de CUDA proviene de la ruta estratégica de ``asumir pérdidas primero a cambio de escala''. GeForce asumió el costo inicial, y una vez formada la base instalada, esta retroalimentó al negocio de centros de datos. Es una ruta típica de ``el ecosistema primero, el hardware cobra después''.
